\documentclass[11pt,a4paper]{article}
\usepackage[utf8]{inputenc}
\usepackage[T1]{fontenc}
\usepackage{lmodern}
\usepackage[french]{babel}
\usepackage{geometry}
\usepackage{hyperref}
\usepackage{enumitem}
\usepackage{listings}
\geometry{margin=2.2cm}
\hypersetup{colorlinks=true,linkcolor=blue,urlcolor=blue}

\lstdefinestyle{cmd}{
  basicstyle=\ttfamily\small,
  breaklines=true,
  frame=single,
  columns=fullflexible,
  keepspaces=true
}

\title{Rapport d'equipe\\Ajout d'une nouvelle serie STM32Cube dans le pipeline}
\author{Projet STM32Cube Preprocessing}
\date{\today}

\begin{document}
\maketitle
\tableofcontents
\newpage

\section{Objet du rapport}
Ce document formalise la procedure standard d'ajout d'une nouvelle serie STM32Cube
(exemple: C5) dans le pipeline. Il precise les fichiers a modifier, la signification
des champs de configuration, la source des donnees a collecter, l'integration dans
l'automatisation, ainsi que les etapes de validation avant mise en production.

Ce rapport est concu pour etre partage avec le document complementaire
\emph{Parametres du workflow Manual Auto Update KB et demarrage rapide du runner
self-hosted} (\texttt{workflow\_dispatch\_inputs\_and\_runner\_quickstart.pdf}), qui
detaille chaque parametre du workflow de declenchement manuel ainsi que
l'installation d'un runner self-hosted de test. Les deux documents forment ensemble
la documentation de reference remise a l'encadrant.

\section{Perimetre technique}
Le perimetre couvre:
\begin{itemize}[leftmargin=*]
  \item la configuration serie,
  \item l'execution du pipeline (Prepare/Upload),
  \item le split des issues par composant,
  \item l'upload vers la KB,
  \item l'integration dans les workflows GitHub Actions,
  \item le test sur une autre machine avec self-hosted runner.
\end{itemize}

\section{Fichiers a modifier}
\subsection{Configuration}
\begin{itemize}[leftmargin=*]
  \item \texttt{shared/config/config\_c5.json} (nouveau fichier)
  \item \texttt{shared/config/config\_all\_series.json} (mise a jour)
\end{itemize}

\subsection{Scripts workflow}
\begin{itemize}[leftmargin=*]
  \item \texttt{pipeline\_Automation/workflow/Run\_Single\_Series\_Full\_Pipeline\_And\_Upload.ps1}
  \item \texttt{pipeline\_Automation/workflow/Split\_Issues\_By\_Component.ps1}
  \item \texttt{pipeline\_Automation/workflow/Upload\_Issues\_By\_Component.ps1}
  \item \texttt{pipeline\_Automation/workflow/run\_all\_series.py}
\end{itemize}

\subsection{Workflows CI/CD}
\begin{itemize}[leftmargin=*]
  \item \texttt{.github/workflows/auto\_update\_kb.yml}
  \item \texttt{.github/workflows/manual\_auto\_update\_kb.yml}
\end{itemize}

\section{Definition des champs de configuration}
\subsection{Fichier \texttt{config\_<serie>.json}}
\textbf{owner}\\
Organisation source des depots (STMicroelectronics).

\textbf{repos}\\
Liste ordonnee des depots de la serie. Le premier element doit etre le repo parent
(exemple: STM32CubeC5), suivi des depots lies (HAL/CMSIS/BSP/Middleware).

\textbf{data\_dir}\\
Dossier des artefacts data intermediaires (valeur standard: \texttt{data}).

\textbf{repos\_dir}\\
Dossier local des clones (valeur standard: \texttt{GitHub\_repos}).

\textbf{auto\_clone\_missing\_repos}\\
Activation du clonage automatique des depots manquants.

\textbf{repo\_clone\_urls}\\
Dictionnaire \texttt{repo -> URL git} pour les depots relies.

\textbf{linked\_repo\_refs}\\
Dictionnaire \texttt{repo -> SHA} correspondant aux refs des sous-modules pour la
release cible. Ce champ est essentiel pour la reproductibilite.

\textbf{commit\_pr\_scope}\\
Strategie de liaison des commits/PR (valeur recommandee: \texttt{linked\_or\_merged}).

\textbf{github\_token\_env}\\
Nom de variable d'environnement contenant le token GitHub.

\textbf{rate\_limits}\\
Temporisations API (issues/comments) pour stabiliser les appels.

\textbf{ingest\_version}\\
Version logique de la configuration d'ingestion.

\subsection{Fichier \texttt{config\_all\_series.json}}
\textbf{all\_series}\\
Liste des series actives.

\textbf{series\_config\_files}\\
Mapping \texttt{Serie -> fichier config}.

\textbf{kb\_datasource\_ids}\\
IDs des datasources par serie (issues, files, diagnostic, resolver).

\textbf{kb\_id}\\
Identifiant de la Knowledge Base cible.

\section{Source des informations: repos, urls, refs}
Les informations sont extraites du repo parent de la serie, positionne sur la
release cible.

\subsection{Commandes de collecte}
\label{sec:commandes-collecte}
\begin{lstlisting}[style=cmd]
git clone --recursive https://github.com/STMicroelectronics/STM32CubeC5.git
cd STM32CubeC5
git checkout <tag_cible>
git submodule update --init --recursive

# URLs des sous-modules
git config -f .gitmodules --get-regexp "submodule\..*\.url"

# SHAs de references
git submodule status --recursive

# Paths des sous-modules
git config -f .gitmodules --get-regexp "submodule\..*\.path"
\end{lstlisting}

\subsection{Regles de remplissage}
\begin{itemize}[leftmargin=*]
  \item \texttt{repos[0]} = repo parent de la serie.
  \item \texttt{repo\_clone\_urls} renseigne les depots a cloner.
  \item \texttt{linked\_repo\_refs} renseigne les SHAs de la release cible.
  \item Exclure les depots non relies au perimetre de la serie.
\end{itemize}

\section{Procedure d'integration d'une nouvelle serie}
\subsection{Etape 1: creation du fichier config serie}
Creer \texttt{shared/config/config\_c5.json} a partir d'un modele proche
(\texttt{config\_c0.json} ou \texttt{config\_g4.json}) et renseigner les champs
\texttt{repos}, \texttt{repo\_clone\_urls}, \texttt{linked\_repo\_refs}.

\subsection{Etape 2: declaration globale}
Mettre a jour \texttt{shared/config/config\_all\_series.json}:
\begin{itemize}[leftmargin=*]
  \item ajout dans \texttt{all\_series},
  \item mapping dans \texttt{series\_config\_files},
  \item entree dans \texttt{kb\_datasource\_ids}.
\end{itemize}

\subsection{Etape 3: mapping scripts}
Ajouter la serie dans tous les \texttt{ValidateSet} et \texttt{seriesMap}:
\begin{itemize}[leftmargin=*]
  \item code serie: \texttt{C5}
  \item slug serie: \texttt{stm32cubec5}
  \item repo parent: \texttt{STM32CubeC5}
  \item config: \texttt{shared/config/config\_c5.json}
\end{itemize}

\subsection{Etape 4: mapping batch Python}
Dans \texttt{run\_all\_series.py}, ajouter \texttt{c5} dans \texttt{SERIES\_CONFIGS}.

\subsection{Etape 5: mapping GitHub Actions}
Ajouter la serie C5 dans:
\begin{itemize}[leftmargin=*]
  \item les listes de series autorisees,
  \item les listes par defaut (si desired),
  \item les normalisations d'inputs.
\end{itemize}

\section{Validation fonctionnelle}
\subsection{Prepare}
\begin{lstlisting}[style=cmd]
powershell -ExecutionPolicy Bypass -File .\pipeline_Automation\workflow\Run_Single_Series_Full_Pipeline_And_Upload.ps1 -Series C5 -Mode Prepare
\end{lstlisting}

\subsection{Split composant}
\begin{lstlisting}[style=cmd]
powershell -ExecutionPolicy Bypass -File .\pipeline_Automation\workflow\Split_Issues_By_Component.ps1 -Series C5
\end{lstlisting}

\subsection{Controle couverture et doublons}
\begin{lstlisting}[style=cmd]
python -m pipeline.delivery.validate_issues_by_component_split --summary "datasets/07_delivery/st_ready/by_series/stm32cubec5/issues_json/by_component/summary_by_component_c5.json"
\end{lstlisting}

\subsection{Upload composant}
\begin{lstlisting}[style=cmd]
powershell -ExecutionPolicy Bypass -File .\pipeline_Automation\workflow\Upload_Issues_By_Component.ps1 -Series C5 -DryRun -RemoteUser "<email_st>"
powershell -ExecutionPolicy Bypass -File .\pipeline_Automation\workflow\Upload_Issues_By_Component.ps1 -Series C5 -RemoteUser "<email_st>"
\end{lstlisting}

\section{Test sur une autre machine (self-hosted runner)}
\subsection{Prerequis}
\begin{itemize}[leftmargin=*]
  \item Windows 10/11
  \item Git
  \item Python 3.12+
  \item acces reseau GitHub/ST
\end{itemize}

\subsection{Installation projet}
\begin{lstlisting}[style=cmd]
git clone https://github.com/<org-ou-user>/PFE_Chatbot_STM32Cube.git
cd PFE_Chatbot_STM32Cube
python -m venv .venv
.\.venv\Scripts\Activate.ps1
python -m pip install --upgrade pip
pip install -r requirements.txt
\end{lstlisting}

\subsection{Runner et configuration}
\begin{itemize}[leftmargin=*]
  \item Declarer un self-hosted runner depuis les settings du repo, en suivant la
  documentation officielle GitHub: \url{https://docs.github.com/en/actions/how-tos/manage-runners/self-hosted-runners/add-runners}.
  \item Lancer le runner en service.
  \item Verifier le label \texttt{self-hosted}.
  \item Activer \texttt{git core.longpaths=true}.
  \item Configurer secrets: \texttt{ST\_GITHUB\_ANALYZER\_API\_KEY}, \texttt{ST\_CHATGPT\_API\_KEY},
  \texttt{ST\_REMOTE\_USER}.
\end{itemize}

\textbf{Important -- clone du depot vs reconfiguration du runner.} Lorsqu'une nouvelle
machine clone ce depot, certaines etapes ont deja ete realisees (scripts,
configuration pipeline, dossier \texttt{actions-runner} present dans l'arborescence
clonee). En revanche, le runner self-hosted lui-meme n'est \textbf{pas}
reutilisable tel quel: il doit etre reconfigure sur la nouvelle machine. Pour cela:
\begin{enumerate}[leftmargin=*]
  \item Supprimer le contenu du dossier \texttt{actions-runner} (le dossier provient
  de la machine d'origine et contient une configuration/registration qui lui est
  propre).
  \item Reprendre les instructions a partir de l'\textbf{etape 6} du document GitHub
  officiel cite ci-dessus (telechargement et configuration du runner).
  \item Executer ces commandes depuis une console \textbf{PowerShell} ou
  \textbf{CMD} Windows, et non depuis Git Bash (le script de configuration du
  runner n'est pas garanti de fonctionner correctement sous Git Bash).
  \item Lors de la (re)configuration (\texttt{config.cmd}), il est recommande de
  conserver les options par defaut proposees (nom de runner, groupe, labels, dossier
  de travail), sauf besoin explicite de les modifier.
\end{enumerate}

\subsection{Test de bout en bout}
\begin{lstlisting}[style=cmd]
powershell -ExecutionPolicy Bypass -File .\pipeline_Automation\workflow\Run_Single_Series_Full_Pipeline_And_Upload.ps1 -Series C5 -Mode Full -EnableIssuesByComponentFlow -ComponentIssuesOnly -RemoteUser "<email_st>"
\end{lstlisting}

\section{Mode migration vers datasources par composant}
Pour remplacer le datasource global issues par des datasources
\texttt{issues\_<SERIE>\_<COMPONENT>}:
\begin{itemize}[leftmargin=*]
  \item activer \texttt{EnableIssuesByComponentFlow},
  \item activer \texttt{ComponentIssuesOnly}.
\end{itemize}

\section{Checklist Go/No-Go}
\begin{itemize}[leftmargin=*]
  \item config serie creee et controlee,
  \item config globale mise a jour,
  \item mappings scripts a jour,
  \item mappings workflows a jour,
  \item phase Prepare validee,
  \item split composant valide,
  \item validation anti-doublons OK,
  \item upload dry-run OK,
  \item upload reel OK,
  \item datasource IDs traces dans \texttt{config\_all\_series.json}.
\end{itemize}

\section{Points d'attention}
Cette section detaille chacun des quatre risques operationnels identifies pour
l'ajout d'une nouvelle serie. Pour chaque point: la notion concernee, pourquoi
c'est un risque, un exemple concret, et comment le detecter/corriger.

\subsection{SHAs incorrects dans \texttt{linked\_repo\_refs}}
\textbf{Notion.} Le champ \texttt{linked\_repo\_refs} du fichier
\texttt{config\_<serie>.json} associe chaque sous-depot (HAL, CMSIS, BSP,
Middleware) au SHA de commit exact correspondant a la release ciblee du depot
parent. C'est ce SHA, et non une branche mobile, qui garantit que le contenu
ingere correspond a une version figee et tracable.

\textbf{Pourquoi c'est un risque.} Si un SHA est incorrect (copie depuis la
mauvaise branche, depot desynchronise, ou oubli de mise a jour apres un
changement de release), le pipeline ingere un etat du sous-depot qui ne
correspond pas a la release officielle. Les documents generes referencent alors
un contenu different de celui reellement publie, ce qui casse la
\emph{reproductibilite}: deux executions du pipeline a des dates differentes,
ou par deux personnes differentes, ne produisent plus le meme resultat pour la
meme release nominale.

\textbf{Exemple concret.} Un SHA laisse sur la branche \texttt{main} au lieu du
tag \texttt{v1.2.0} de la release: si le sous-depot recoit de nouveaux commits
apres l'integration, la prochaine execution du pipeline changera silencieusement
les donnees ingerees, sans que la version de la serie declaree ait change.

\textbf{Detection et correction.} Toujours extraire les SHAs avec
\texttt{git submodule status --recursive} apres avoir positionne le depot parent
exactement sur le tag/commit de la release cible (voir section
\ref{sec:commandes-collecte} sur les commandes de collecte), jamais a la main
ni depuis une branche en mouvement.

\subsection{Serie declaree en config mais absente des \texttt{ValidateSet}/\texttt{seriesMap}}
\textbf{Notion.} \texttt{ValidateSet} est une contrainte PowerShell qui restreint
les valeurs acceptees pour un parametre de script (ex: liste figee de codes
serie). \texttt{seriesMap} est la table interne (PowerShell ou Python) qui relie
un code serie (ex: \texttt{C5}) a son fichier de configuration, son slug et son
depot parent.

\textbf{Pourquoi c'est un risque.} Declarer une serie uniquement dans
\texttt{config\_all\_series.json} ne suffit pas si un script utilise encore une
liste \texttt{ValidateSet} ou une table \texttt{seriesMap} codee en dur: le
script refuse la serie ou echoue en lui cherchant un fichier de configuration
inexistant, meme si la configuration JSON globale est correcte.

\textbf{Etat verifie pour la serie C5 (mise a jour de ce rapport).} Une
verification du code a ete effectuee. Les scripts principaux
\texttt{Run\_Single\_Series\_Full\_Pipeline\_And\_Upload.ps1},
\texttt{Split\_Issues\_By\_Component.ps1}, \texttt{Upload\_Issues\_By\_Component.ps1}
et \texttt{Upload\_Series\_All\_Categories.ps1} construisent leur
\texttt{seriesMap} \emph{dynamiquement} a partir de
\texttt{shared/config/config\_all\_series.json}: comme \texttt{STM32CubeC5} et
\texttt{config\_c5.json} y sont bien enregistres, \textbf{ce risque est corrige
pour ces quatre scripts et la serie C5 est utilisable directement}.

Toutefois, deux scripts auxiliaires plus anciens utilisent encore des tables
\emph{statiques} qui n'incluent pas C5 et devront etre mis a jour manuellement
avant utilisation avec cette serie:
\begin{itemize}[leftmargin=*]
  \item \texttt{pipeline\_Automation/workflow/run\_all\_series.py} (dictionnaire
  \texttt{SERIES\_CONFIGS} code en dur, C5 absent),
  \item \texttt{pipeline\_Automation/workflow/reexport\_delivery.py} (meme
  dictionnaire \texttt{SERIES\_CONFIGS} statique, C5 absent),
  \item \texttt{pipeline\_Automation/workflow/Upload\_Multi\_Series\_All\_Categories.ps1}
  (liste \texttt{ValidateSet} figee, C5 absent).
\end{itemize}
Ces trois scripts restent \textbf{a risque} tant qu'ils ne sont pas mis a jour:
leur usage avec \texttt{-Series C5} (ou \texttt{--series c5}) echouera avec une
erreur de type \og config introuvable \fg{} ou \og valeur non autorisee \fg{}.

\textbf{Detection et correction.} Avant d'utiliser un script pour une serie
nouvellement ajoutee, verifier s'il lit \texttt{config\_all\_series.json}
dynamiquement (pas de risque) ou s'il possede une liste/dictionnaire de series
code en dur (risque residuel, necessite un ajout manuel de la serie dans ce
fichier).

\subsection{Serie absente des workflows}
\textbf{Notion.} Les workflows GitHub Actions
(\texttt{.github/workflows/auto\_update\_kb.yml} et
\texttt{.github/workflows/manual\_auto\_update\_kb.yml}) exposent les series
disponibles via leurs parametres \texttt{workflow\_dispatch} (voir le document
complementaire sur les parametres du workflow). Meme si une serie est
correctement configuree cote pipeline, elle doit aussi etre acceptee par la
logique de normalisation des inputs du workflow.

\textbf{Pourquoi c'est un risque.} Si la serie n'est pas prise en compte dans
cette normalisation (ou dans une liste par defaut du workflow), il devient
impossible de la declencher depuis l'interface \emph{Actions} de GitHub, meme en
executant le script localement avec succes: le workflow rejettera l'entree ou
ne l'inclura pas dans la matrice de jobs.

\textbf{Etat verifie pour C5.} Les workflows lisent egalement
\texttt{config\_all\_series.json} pour deriver la liste des series autorisees
(meme mecanisme que \texttt{seriesMap}), donc C5 y est normalement deja
acceptee. Il reste toutefois recommande de faire un essai de declenchement
manuel avec \texttt{series=C5} et \texttt{mode=Prepare} pour confirmer avant un
usage en production.

\textbf{Detection et correction.} Lancer le workflow \emph{Manual Auto Update KB}
avec la nouvelle serie en mode \texttt{Prepare} et verifier que le job
correspondant apparait dans la matrice, avant de tenter un mode \texttt{Full}.

\subsection{Upload composant sans split prealable}
\textbf{Notion.} Le \og split par composant \fg{} (script
\texttt{Split\_Issues\_By\_Component.ps1}) decoupe le datasource \og issues \fg{}
monolithique d'une serie en plusieurs datasources plus petits, un par composant
logiciel (HAL, BSP, Middleware, etc.), et genere en meme temps un fichier
\texttt{summary\_by\_component\_<serie>.json} qui recapitule la repartition et
sert de reference pour le controle anti-doublons.

\textbf{Pourquoi c'est un risque.} Le script
\texttt{Upload\_Issues\_By\_Component.ps1} attend en entree ce fichier
\texttt{summary\_by\_component} pour savoir quels composants uploader et avec
quels identifiants de datasource. Si l'upload par composant est lance sans avoir
prealablement execute le split, ce fichier de synthese est absent et l'upload ne
peut pas determiner la liste des composants a traiter: il echoue immediatement
ou traite une liste vide.

\textbf{Exemple concret.} Executer directement
\texttt{Upload\_Issues\_By\_Component.ps1 -Series C5} sans avoir lance au
prealable \texttt{Split\_Issues\_By\_Component.ps1 -Series C5} produit une erreur
de fichier manquant plutot qu'un upload partiel silencieux.

\textbf{Detection et correction.} Toujours respecter l'ordre: (1) export
delivery standard, (2) \texttt{Split\_Issues\_By\_Component.ps1}, (3) controle du
resume via \texttt{validate\_issues\_by\_component\_split}, (4) uniquement
ensuite \texttt{Upload\_Issues\_By\_Component.ps1}. Cet ordre est deja formalise
dans la section \og Validation fonctionnelle \fg{} de ce document.

\end{document}
